%! TEX root = ../main.tex
\documentclass[../main.tex]{subfiles}

\begin{document}

\begin{abstract}
    Flying insect-scale micro-robots have been in development since early 2000\cite{4457871}. These types of robots have demonstrated exceptional maneuverability with insect-like flight capabilities\cite{9357346} at only a fraction of the weight and size of the lightest commercial quadcopter\cite{8046794}. These characteristics are essential for navigating small nooks and crannies, potentially allowing these robots to survey and monitor areas that are otherwise inaccessible to human and larger robots. However, due to the limited payload associated with small-scale systems, these robots are shackled to laboratory environments consist of motion capture systems, powerful computers, and external high-voltage amplifiers for sensing, controlling, and powering\cite{Hsiao_Kim_Ceron_Ren_Chen_2023}. This research project aims to untether these robots from the lab bench by developing light-weight sensing, control, and power electronics systems that these robots can carry for autonomous operations.
    
    % However, current implementations of these robots relies on piezoelectric actuators that are hard to scale and pose a significant challenge to miniaturization \cite{}. Therefore, 
    
    % power electronics circuits for autonomous
\end{abstract}

\end{document}


